
\section*{7. Appendice : Spécialisation en théorie des intersections}
\begin{center}
\underline{par A. Grothendieck}
\end{center}

\textbf{Lemme 7.1.} \uline{Soit} @@M0@@ \uline{une immersion fermée régulière} (VII 1.4)\uline{, de codimension bornée, de sorte que} @@M1@@ \uline{est un morphisme parfait} (VII 1.9)\uline{, et qu'on peut donc définir le morphisme} (IV 2.12)
@@M2@@
\uline{Soit d'autre part} @@M3@@ \uline{le faisceau conormal pour} @@M4@@\uline{, qui est un faisceau localement libre de type fini sur} @@M5@@ (VII 1.6)\uline{, de sorte qu'on peut former}
@@M6@@
\uline{Considérons aussi l'homomorphisme image directe} (IV 2.11)
@@M7@@
\uline{Ceci posé, on a la formule}
@@M8@@
\uline{où le produit du deuxième membre est relatif à la structure de} @@M9@@\uline{-module de} @@M10@@ \uline{définie dans} (IV 2.10).

La démonstration est, essentiellement sans changement, la même que pour la même formule, écrite en travaillant avec @@M11@@, @@M12@@, établie dans (VII 2.7).

\textbf{Corollaire 7.2.} \uline{Soient} @@M13@@ \uline{un morphisme plat, avec} @@M14@@ \uline{un trait i.e. le spectre d'un anneau de valuation discrète} @@M15@@, @@M16@@ \uline{le point fermé de} @@M17@@, @@M18@@ \uline{son point générique,} @@M19@@ \uline{l'inclusion canonique. Alors} @@M20@@ \uline{est de tor-dimension finie donc} (IV 2.12) \uline{on peut définir}
@@M21@@
